\documentclass[12pt,a4paper]{article}
\usepackage[utf8]{inputenc}
\usepackage[indonesian]{babel}
\usepackage{geometry}
\geometry{top=2.5cm, bottom=2.5cm, left=3cm, right=2.5cm}
\usepackage{hyperref}
\usepackage{xcolor}
\usepackage{listings}
\usepackage{tcolorbox}
\usepackage{booktabs}
\usepackage{tabularx}
\usepackage{enumitem}
\usepackage{titlesec}
\usepackage{tikz}
\usetikzlibrary{shapes.geometric, arrows.meta, positioning, shadows, calc}

\hypersetup{
    colorlinks=true,
    linkcolor=teal,
    filecolor=magenta,      
    urlcolor=blue,
    pdftitle={Panduan Presentasi Project SpotKu},
}

% Pengaturan Warna Code Snippet & TikZ
\definecolor{codegreen}{rgb}{0,0.6,0}
\definecolor{codegray}{rgb}{0.5,0.5,0.5}
\definecolor{codepurple}{rgb}{0.58,0,0.82}
\definecolor{backcolour}{rgb}{0.96,0.96,0.96}
\definecolor{boxblue}{rgb}{0.90,0.94,0.98}
\definecolor{darkblue}{rgb}{0.10,0.30,0.60}
\definecolor{nodeblue}{rgb}{0.85,0.92,1.0}
\definecolor{nodegreen}{rgb}{0.85,0.95,0.88}
\definecolor{nodeorange}{rgb}{1.0,0.92,0.85}

% Gaya Node TikZ
\tikzstyle{startstop} = [rectangle, rounded corners=12pt, minimum width=3cm, minimum height=1cm, text centered, draw=darkblue, fill=nodeblue, font=\bfseries\small, drop shadow]
\tikzstyle{screen} = [rectangle, rounded corners=6pt, minimum width=3.2cm, minimum height=1cm, text centered, draw=teal!80!black, fill=nodegreen, text width=3cm, font=\small, drop shadow]
\tikzstyle{process} = [rectangle, rounded corners=6pt, minimum width=3.2cm, minimum height=1cm, text centered, draw=blue!80!black, fill=boxblue, text width=3cm, font=\small, drop shadow]
\tikzstyle{decision} = [diamond, aspect=2, minimum width=2.8cm, minimum height=0.8cm, text centered, draw=orange!90!black, fill=nodeorange, text width=2.6cm, font=\scriptsize\bfseries, drop shadow]
\tikzstyle{arrow} = [thick,->,>=stealth,color=darkblue]

\lstdefinestyle{mystyle}{
    backgroundcolor=\color{backcolour},   
    commentstyle=\color{codegreen},
    keywordstyle=\color{blue}\bfseries,
    numberstyle=\tiny\color{codegray},
    stringstyle=\color{codepurple},
    basicstyle=\ttfamily\footnotesize,
    breakatwhitespace=false,         
    breaklines=true,                 
    captionpos=b,                    
    keepspaces=true,                 
    numbers=left,                    
    numbersep=5pt,                  
    showspaces=false,                
    showstringspaces=false,
    showtabs=false,                  
    tabsize=2
}
\lstset{style=mystyle}

\tcbuselibrary{skins,breakable}
\newtcolorbox{highlightbox}[1][]{
  colback=boxblue,
  colframe=darkblue,
  fonttitle=\bfseries,
  title=#1,
  arc=2mm,
  breakable
}

\title{
    \textbf{\Large BUKU PANDUAN \& DOKUMENTASI TEKNIS PRESENTASI PROJECT SPOTKU}\\[0.5em]
    \large Perancangan Aplikasi Diari Perjalanan \& Kenangan Berbasis Mobile (Android)\\
    menggunakan Flutter, Firebase, dan OpenStreetMap
}
\author{\textbf{Dokumentasi Akademik Presentasi Dosen}}
\date{\today}

\begin{document}

\maketitle
\tableofcontents
\newpage

\section{Pendahuluan \& Gambaran Umum (Elevator Pitch)}

\textbf{SpotKu} adalah aplikasi \textit{mobile} berbasis Android yang dirancang sebagai diari perjalanan digital (\textit{Travel \& Memory Diary}). Aplikasi ini memungkinkan pengguna mencatat tempat-tempat berkesan yang dikunjungi lengkap dengan foto/video lokasi, catatan pribadi, serta penandaan koordinat geografis presisi di atas peta interaktif.

\subsection{Spesifikasi Teknologi (Tech Stack)}
\begin{itemize}[leftmargin=*]
    \item \textbf{Bahasa Pemrograman}: Dart (SDK $\ge 3.3.0$)
    \item \textbf{Framework UI}: Flutter (Material Design 3)
    \item \textbf{Autentikasi}: Firebase Auth (Google Sign-In via OAuth 2.0)
    \item \textbf{Database Cloud}: Cloud Firestore (NoSQL Realtime Database)
    \item \textbf{Layanan Peta}: OpenStreetMap via \texttt{flutter\_map} (Tanpa Google Maps Billing)
    \item \textbf{Integrasi Hardware HP}: 
    \begin{itemize}
        \item \texttt{geolocator}: Membaca posisi GPS (Latitude \& Longitude)
        \item \texttt{image\_picker}: Mengakses Kamera HP untuk foto dan video
        \item \texttt{flutter\_local\_notifications}: Memicu notifikasi lokal Android
    \end{itemize}
\end{itemize}

---

\section{Use Case \& Workflow Visual (TikZ Diagram)}

\subsection{Daftar Use Case Sistem}

\begin{table}[h!]
\centering
\small
\begin{tabularx}{\textwidth}{l l l X l}
\toprule
\textbf{ID} & \textbf{Nama Use Case} & \textbf{Aktor} & \textbf{Deskripsi Singkat} & \textbf{Layanan/Hardware} \\
\midrule
UC-01 & Login Akun Google & User & Masuk ke aplikasi menggunakan akun Gmail tanpa password manual. & Firebase Auth \& OAuth 2.0 \\
UC-02 & Peta Spot Realtime & User & Melihat peta dunia dengan pin merah lokasi kenangan. & OpenStreetMap \& Firestore \\
UC-03 & Tambah Spot Baru & User & Memotret foto/video, baca GPS otomatis, dan simpan catatan. & Kamera, GPS, \& Firestore \\
UC-04 & Lihat Detail Spot & User & Mengklik pin merah di peta untuk melihat foto \& koordinat. & Memori Lokal HP \\
UC-05 & Notifikasi Pengingat & User & Menjalankan timer 5 detik untuk menguji notifikasi Android. & Local Notifications \\
UC-06 & Logout & User & Keluar dari sesi akun dan kembali ke layar login. & Firebase Auth \\
\bottomrule
\end{tabularx}
\caption{Daftar Spesifikasi Use Case Aplikasi SpotKu}
\end{table}

\newpage
\subsection{Diagram Vektor Visual Alur Kerja Pengguna (TikZ Flowchart)}

Diagram di bawah ini secara visual memperlihatkan bagaimana pengguna berinteraksi dengan sistem aplikasi SpotKu dari proses awal pembukaan aplikasi hingga penambahan tempat baru:

\begin{figure}[h!]
\centering
\begin{tikzpicture}[node distance=1.8cm and 1.5cm, auto]
    % Nodes
    \node (start) [startstop] {START: Buka App};
    \node (auth) [decision, below=of start] {Status Login Firebase?};
    \node (login) [screen, left=of auth] {LoginScreen\\(Google Sign-In)};
    \node (oauth) [process, below=of login] {Otentikasi OAuth 2.0\\Firebase Auth};
    \node (home) [screen, below=of auth, yshift=-0.5cm] {HomeScreen\\(Peta OpenStreetMap)};
    \node (action) [decision, below=of home] {Pilihan Aksi User?};
    \node (add) [screen, left=of action] {AddSpotScreen\\(Form Kenangan)};
    \node (gpscam) [process, below=of add] {1. GPS (Geolocator)\\2. Kamera (ImagePicker)};
    \node (firestore) [process, below=of gpscam] {Simpan Dokumen ke\\Cloud Firestore};
    \node (detail) [screen, right=of action] {DetailScreen\\(Detail Spot \& Foto)};
    \node (notif) [process, below=of detail] {Notifikasi Timer 5s\\(Local Notifications)};

    % Lines / Arrows
    \draw [arrow] (start) -- (auth);
    \draw [arrow] (auth) -- node[anchor=south] {Belum} (login);
    \draw [arrow] (login) -- (oauth);
    \draw [arrow] (oauth) |- (home);
    \draw [arrow] (auth) -- node[anchor=west] {Sudah} (home);
    \draw [arrow] (home) -- (action);
    \draw [arrow] (action) -- node[anchor=south] {Klik (+)} (add);
    \draw [arrow] (add) -- (gpscam);
    \draw [arrow] (gpscam) -- (firestore);
    \draw [arrow] (firestore) |- (home);
    \draw [arrow] (action) -- node[anchor=south] {Klik Pin} (detail);
    \draw [arrow] (detail) -- (notif);
\end{tikzpicture}
\caption{Diagram Visual Alur Kerja Sistem SpotKu (TikZ Vector Flowchart)}
\end{figure}

---

\section{Rasionalitas Keputusan Desain (Design Rationale)}

\begin{highlightbox}[Mengapa Menggunakan OpenStreetMap daripada Google Maps SDK?]
\textbf{Alasan Utama}: Google Maps SDK memerlukan pendaftaran kartu kredit (\textit{billing account}) di Google Cloud Platform untuk mendapatkan API Key. Jika limit gratis terlampaui, Google Maps akan eror atau menagih biaya.\\
\textbf{Keunggulan}: OpenStreetMap bersifat 100\% \textit{open-source} dan gratis tanpa API Key. Paket \texttt{flutter\_map} mengunduh ubin peta (\textit{map tiles}) secara langsung dan sangat ringan di Android.
\end{highlightbox}

\begin{highlightbox}[Mengapa Berkas Foto Disimpan secara Lokal (\texttt{media\_url})?]
\textbf{Alasan Utama}: Mengunggah file media berukuran besar ke Firebase Storage membutuhkan koneksi internet yang kencang, kuota besar, dan biaya langganan cloud storage.\\
\textbf{Keunggulan}: Dengan menyimpan alamat path file di memori internal HP (misal: \texttt{/data/user/0/.../image.jpg}), proses simpan spot dapat dilakukan secara \textbf{seketika (\textit{instant})} tanpa bergantung pada kecepatan upload internet.
\end{highlightbox}

\begin{highlightbox}[Mengapa Menggunakan \texttt{StreamBuilder} di Halaman Utama?]
\textbf{Alasan Utama}: \texttt{FutureBuilder} hanya membaca data satu kali saat halaman pertama dibangun.\\
\textbf{Keunggulan}: \texttt{StreamBuilder} membuka koneksi \textit{WebSocket} secara persisten dengan Cloud Firestore. Ketika ada data spot baru ditambahkan, marker di peta langsung diperbarui secara \textit{realtime} tanpa perlu melakukan \textit{refresh} manual.
\end{highlightbox}

---

\section{Pembedahan Kode Berbaris \& Pengaruhnya ke Aplikasi}

\subsection{File 1: Konfigurasi \texttt{pubspec.yaml}}

\begin{lstlisting}[language=XML, firstnumber=9, caption=pubspec.yaml]
dependencies:
  flutter:
    sdk: flutter
  cupertino_icons: ^1.0.8

  # Autentikasi (gratis - Firebase Spark plan)
  firebase_core: ^3.6.0
  firebase_auth: ^5.3.1
  google_sign_in: ^6.2.1

  # Database (gratis - Firestore free tier)
  cloud_firestore: ^5.4.3

  # Kamera (gratis, tanpa API key)
  image_picker: ^1.1.2

  # Peta (GRATIS, OpenStreetMap, tanpa API key / billing)
  flutter_map: ^7.0.2
  latlong2: ^0.9.1

  # Lokasi GPS (gratis)
  geolocator: ^13.0.1

  # Notifikasi lokal (gratis)
  flutter_local_notifications: ^18.0.1
  intl: ^0.19.0
\end{lstlisting}

\begin{highlightbox}[Penjelasan Rinci Nomor Baris \texttt{pubspec.yaml}]
\begin{itemize}[leftmargin=*]
    \item \textbf{Baris 15--17}: Paket dasar Firebase, pengelola sesi login, dan pustaka penyedia dialog Google Sign-In.
    \item \textbf{Baris 20}: Pustaka basis data cloud NoSQL Firestore untuk menyimpan dokumen spot kenangan.
    \item \textbf{Baris 23}: Pustaka penghubung antarmuka aplikasi dengan kamera perangkat Android.
    \item \textbf{Baris 26--27}: Pustaka perender peta ubin OpenStreetMap (\texttt{flutter\_map}) dan kalkulator objek geografis (\texttt{latlong2}).
    \item \textbf{Baris 30}: Pustaka pembaca chip sensor GPS perangkat (\texttt{geolocator}).
    \item \textbf{Baris 33}: Pustaka pembuat notifikasi pop-up lokal pada Android.
\end{itemize}
\end{highlightbox}

\newpage
\subsection{File 2: Inisialisasi \& Routing \texttt{lib/main.dart}}

\begin{lstlisting}[language=Java, firstnumber=10, caption=lib/main.dart (main \& AuthGate)]
Future<void> main() async {
  WidgetsFlutterBinding.ensureInitialized();
  
  try {
    await Firebase.initializeApp(
      options: DefaultFirebaseOptions.currentPlatform,
    );
    await NotificationService.init();
  } catch (e) {
    debugPrint("Error saat inisialisasi Firebase/Notifikasi: $e");
  }
  
  runApp(const SpotKuApp());
}

class AuthGate extends StatelessWidget {
  const AuthGate({super.key});

  @override
  Widget build(BuildContext context) {
    return StreamBuilder<User?>(
      stream: FirebaseAuth.instance.authStateChanges(),
      builder: (context, snapshot) {
        if (snapshot.connectionState == ConnectionState.waiting) {
          return const Scaffold(
            body: Center(child: CircularProgressIndicator()),
          );
        }
        if (snapshot.hasData) {
          return const HomeScreen();
        }
        return const LoginScreen();
      },
    );
  }
}
\end{lstlisting}

\begin{highlightbox}[Penjelasan Rinci Nomor Baris \texttt{lib/main.dart}]
\begin{itemize}[leftmargin=*]
    \item \textbf{Baris 11 (\texttt{WidgetsFlutterBinding.ensureInitialized()})}:\\
    \textit{Pengaruh}: Memastikan sistem native Android telah siap menerima pemanggilan kode asinkron sebelum fungsi \texttt{runApp} dijalankan.
    \item \textbf{Baris 14--17 (\texttt{Firebase.initializeApp()} \& \texttt{NotificationService.init()})}:\\
    \textit{Pengaruh}: Menginisialisasi koneksi dengan server Firebase dan mengaktifkan saluran notifikasi Android.
    \item \textbf{Baris 52--54 (\texttt{StreamBuilder<User?>})}:\\
    \textit{Pengaruh}: Mendengarkan status login pengguna secara \textit{realtime}.
    \item \textbf{Baris 60--63 (\texttt{snapshot.hasData})}:\\
    \textit{Pengaruh}: Jika pengguna sudah terautentikasi, layar langsung menampilkan \texttt{HomeScreen}. Jika belum, aplikasi menampilkan \texttt{LoginScreen}.
\end{itemize}
\end{highlightbox}

\newpage
\subsection{File 3: Model Data \texttt{lib/models/spot\_model.dart}}

\begin{lstlisting}[language=Java, firstnumber=3, caption=lib/models/spot\_model.dart]
class Spot {
  final String id;
  final String userId;
  final String title;
  final String description;
  final double latitude;
  final double longitude;
  final String mediaType; // "image" atau "video"
  final String mediaUrl; // jalur file lokal di HP
  final Timestamp? createdAt;

  Spot({
    required this.id,
    required this.userId,
    required this.title,
    required this.description,
    required this.latitude,
    required this.longitude,
    required this.mediaType,
    required this.mediaUrl,
    this.createdAt,
  });

  factory Spot.fromFirestore(DocumentSnapshot doc) {
    final data = doc.data() as Map<String, dynamic>;
    return Spot(
      id: doc.id,
      userId: data['user_id'] ?? '',
      title: data['title'] ?? '',
      description: data['description'] ?? '',
      latitude: (data['latitude'] ?? 0).toDouble(),
      longitude: (data['longitude'] ?? 0).toDouble(),
      mediaType: data['media_type'] ?? 'image',
      mediaUrl: data['media_url'] ?? '',
      createdAt: data['created_at'],
    );
  }
}
\end{lstlisting}

\begin{highlightbox}[Penjelasan Rinci Nomor Baris \texttt{lib/models/spot\_model.dart}]
\begin{itemize}[leftmargin=*]
    \item \textbf{Baris 4--12}: Mendefinisikan atribut data lokasi kenangan (ID, ID User, Judul, Deskripsi, Latitude, Longitude, Tipe Media, Path File, dan Timestamp).
    \item \textbf{Baris 26--39 (\texttt{Spot.fromFirestore})}:\\
    \textit{Pengaruh}: Mengubah dokumen JSON dari Firestore menjadi objek Dart \texttt{Spot}. Operator \texttt{??} memberikan nilai \textit{fallback} default untuk mencegah timbulnya \textit{NullPointer Error}.
\end{itemize}
\end{highlightbox}

\newpage
\subsection{File 4: Layanan Notifikasi \texttt{lib/services/notification\_service.dart}}

\begin{lstlisting}[language=Java, firstnumber=7, caption=lib/services/notification\_service.dart]
  static Future<void> init() async {
    const androidSettings =
        AndroidInitializationSettings('@mipmap/ic_launcher');
    const initSettings = InitializationSettings(android: androidSettings);
    await _plugin.initialize(initSettings);

    // Meminta izin runtime notifikasi di Android 13+ (API 33 ke atas)
    final androidPlugin = _plugin.resolvePlatformSpecificImplementation<
        AndroidFlutterLocalNotificationsPlugin>();
    if (androidPlugin != null) {
      await androidPlugin.requestNotificationsPermission();
    }
  }

  static Future<void> showGeofenceReminder({
    required String placeName,
    required String formattedDate,
  }) async {
    _notificationIdCounter++;
    final currentId = _notificationIdCounter;

    const androidDetails = AndroidNotificationDetails(
      'spotku_geofence_channel',
      'SpotKu Geofence Reminders',
      channelDescription: 'Notifikasi otomatis saat dekat titik kenangan',
      importance: Importance.max,
      priority: Priority.high,
      playSound: true,
      enableVibration: true,
      groupKey: 'com.project.spotkuapp.GEOFENCE_STACK',
    );
    const details = NotificationDetails(android: androidDetails);

    await _plugin.show(
      currentId,
      'SpotKu - ' + placeName,
      'Anda pernah mengunjungi tempat ini. (Ditandai: ' + formattedDate + ')',
      details,
    );
  }
\end{lstlisting}

\begin{highlightbox}[Penjelasan Rinci Nomor Baris \texttt{lib/services/notification\_service.dart}]
\begin{itemize}[leftmargin=*]
    \item \textbf{Baris 7--18 (\texttt{init()})}: Mendaftarkan ikon aplikasi (\texttt{ic\_launcher}) dan memanggil \texttt{requestNotificationsPermission()} untuk meminta izin notifikasi di layar HP Android 13+ (seperti Samsung A14).
    \item \textbf{Baris 20--47 (\texttt{showGeofenceReminder})}: Mengirim notifikasi berulang dengan ID unik dan \texttt{groupKey} agar **menumpuk (*stacking*) layaknya pesan WhatsApp** saat pengguna berada di radius $\le$ 5 meter.
    \item \textbf{Baris 49--51 (\texttt{cancelAllNotifications})}: Membersihkan seluruh tumpukan notifikasi ketika pengguna bergerak keluar dari radius 5 meter.
\end{itemize}
\end{highlightbox}

\newpage
\subsection{File 5: Layar Login \texttt{lib/screens/login\_screen.dart}}

\begin{lstlisting}[language=Java, firstnumber=15, caption=lib/screens/login\_screen.dart]
  Future<void> _signInWithGoogle() async {
    setState(() => _loading = true);
    try {
      final googleUser = await GoogleSignIn().signIn();
      if (googleUser == null) {
        setState(() => _loading = false);
        return;
      }

      final googleAuth = await googleUser.authentication;
      final credential = GoogleAuthProvider.credential(
        accessToken: googleAuth.accessToken,
        idToken: googleAuth.idToken,
      );

      await FirebaseAuth.instance.signInWithCredential(credential);
    } catch (e) {
      if (mounted) {
        ScaffoldMessenger.of(context).showSnackBar(
          SnackBar(content: Text('Gagal masuk: $e')),
        );
      }
    } finally {
      if (mounted) setState(() => _loading = false);
    }
  }
\end{lstlisting}

\begin{highlightbox}[Penjelasan Rinci Nomor Baris \texttt{lib/screens/login\_screen.dart}]
\begin{itemize}[leftmargin=*]
    \item \textbf{Baris 18--22 (\texttt{GoogleSignIn().signIn()})}: Membuka dialog pilihan akun Google bawaan sistem Android.
    \item \textbf{Baris 24--28 (\texttt{GoogleAuthProvider.credential})}: Mengambil token OAuth 2.0 (\texttt{idToken} dan \texttt{accessToken}).
    \item \textbf{Baris 30 (\texttt{signInWithCredential})}: Mengirimkan token ke Firebase Auth untuk memverifikasi autentikasi pengguna.
\end{itemize}
\end{highlightbox}

\newpage
\subsection{File 6: Layar Utama Peta \texttt{lib/screens/home\_screen.dart}}

\begin{lstlisting}[language=Java, firstnumber=34, caption=lib/screens/home\_screen.dart]
      body: StreamBuilder<QuerySnapshot>(
        stream: FirebaseFirestore.instance
            .collection('spots')
            .where('user_id', isEqualTo: uid)
            .snapshots(),
        builder: (context, snapshot) {
          if (snapshot.hasError) {
            return Center(child: Text('Terjadi error: ${snapshot.error}'));
          }
          if (!snapshot.hasData) {
            return const Center(child: CircularProgressIndicator());
          }

          final spots =
              snapshot.data!.docs.map((d) => Spot.fromFirestore(d)).toList();

          final markers = spots.map((spot) {
            return Marker(
              point: LatLng(spot.latitude, spot.longitude),
              width: 44,
              height: 44,
              child: GestureDetector(
                onTap: () {
                  Navigator.push(
                    context,
                    MaterialPageRoute(
                      builder: (_) => DetailScreen(spot: spot),
                    ),
                  );
                },
                child: const Icon(
                  Icons.location_pin,
                  color: Colors.redAccent,
                  size: 40,
                ),
              ),
            );
          }).toList();

          return FlutterMap(
            options: const MapOptions(
              initialCenter: LatLng(-6.8871, 109.7745),
              initialZoom: 13,
            ),
            children: [
              TileLayer(
                urlTemplate: 'https://tile.openstreetmap.org/{z}/{x}/{y}.png',
                userAgentPackageName: 'com.example.spotku',
              ),
              MarkerLayer(markers: markers),
            ],
          );
        },
      ),
\end{lstlisting}

\begin{highlightbox}[Penjelasan Rinci Nomor Baris \texttt{lib/screens/home\_screen.dart}]
\begin{itemize}[leftmargin=*]
    \item \textbf{Baris 35--38 (\texttt{where('user_id', isEqualTo: uid)})}: Menyaring dokumen Firestore agar hanya mengembalikan spot milik pengguna yang sedang aktif login.
    \item \textbf{Baris 50--71 (\texttt{Marker \& GestureDetector})}: Mengkonversi pasangan nilai Latitude dan Longitude menjadi widget Pin Merah yang dapat di-klik untuk berpindah ke \texttt{DetailScreen}.
    \item \textbf{Baris 73--85 (\texttt{FlutterMap \& TileLayer})}: Merender peta OpenStreetMap dengan titik awal Pemalang (-6.8871, 109.7745).
\end{itemize}
\end{highlightbox}

\newpage
\subsection{File 7: Layar Form Tambah Spot \texttt{lib/screens/add\_spot\_screen.dart}}

\begin{lstlisting}[language=Java, firstnumber=32, caption=lib/screens/add\_spot\_screen.dart]
  Future<void> _getCurrentLocation() async {
    setState(() => _gettingLocation = true);
    try {
      final serviceEnabled = await Geolocator.isLocationServiceEnabled();
      if (!serviceEnabled) throw 'Layanan lokasi (GPS) tidak aktif';

      var permission = await Geolocator.checkPermission();
      if (permission == LocationPermission.denied) {
        permission = await Geolocator.requestPermission();
      }

      final pos = await Geolocator.getCurrentPosition();
      setState(() => _position = pos);
    } catch (e) {
      // Handling Error
    }
  }

  Future<void> _pickPhoto() async {
    final file =
        await _picker.pickImage(source: ImageSource.camera, imageQuality: 70);
    if (file != null) {
      setState(() {
        _mediaFile = file;
        _mediaType = 'image';
      });
    }
  }

  Future<void> _save() async {
    final uid = FirebaseAuth.instance.currentUser!.uid;
    await FirebaseFirestore.instance.collection('spots').add({
      'user_id': uid,
      'title': _titleController.text.trim(),
      'description': _descController.text.trim(),
      'latitude': _position!.latitude,
      'longitude': _position!.longitude,
      'media_type': _mediaType,
      'media_url': _mediaFile!.path,
      'created_at': FieldValue.serverTimestamp(),
    });

    if (mounted) Navigator.pop(context);
  }
\end{lstlisting}

\begin{highlightbox}[Penjelasan Rinci Nomor Baris \texttt{lib/screens/add\_spot\_screen.dart}]
\begin{itemize}[leftmargin=*]
    \item \textbf{Baris 35--47 (\texttt{Geolocator.getCurrentPosition()})}: Memeriksa izin lokasi Android dan mengambil koordinat presisi dari sensor GPS HP.
    \item \textbf{Baris 61--62 (\texttt{pickImage})}: Membuka aplikasi kamera untuk memotret lokasi dengan tingkat kompresi 70\%.
    \item \textbf{Baris 104--113 (\texttt{collection('spots').add})}: Menyimpan dokumen baru ke Cloud Firestore.
\end{itemize}
\end{highlightbox}

---

\section{Panduan Step-by-Step Buka File \& Tampilan Aplikasi}

\subsection{Langkah 1: Konfigurasi \& Bahan Pustaka (\texttt{pubspec.yaml})}
\begin{enumerate}[leftmargin=*]
    \item \textbf{Navigasi}: Buka folder \texttt{spotku/} $\rightarrow$ Buka file \texttt{pubspec.yaml} (Baris 15--34).
    \item \textbf{Fungsi \& Logika}: Mendaftarkan paket Firebase, Kamera, GPS, Peta OpenStreetMap, dan Notifikasi. Tanpa file ini, pemanggilan paket akan dianggap sebagai teks asing (\textit{Unresolved identifier}).
    \item \textbf{Tampilan Aplikasi}: Tidak berwujud UI langsung, melainkan mengaktifkan fitur hardware kamera, GPS, dan unduh peta dari internet.
\end{enumerate}

\subsection{Langkah 2: Booting \& Pintu Masuk (\texttt{lib/main.dart})}
\begin{enumerate}[leftmargin=*]
    \item \textbf{Navigasi}: Buka folder \texttt{lib/} $\rightarrow$ Buka file \texttt{lib/main.dart} (Baris 10--26 dan 47--66).
    \item \textbf{Fungsi \& Logika}: Menyiapkan inisialisasi Firebase dan mengecek status login user (\texttt{AuthGate}). \texttt{AuthGate} WAJIB ada agar pengguna yang pernah login tidak perlu dipaksa login berulang kali.
    \item \textbf{Tampilan Aplikasi}: Saat aplikasi dibuka, muncul indikator loading sebentar, lalu aplikasi menentukan apakah membuka \texttt{LoginScreen} atau \texttt{HomeScreen}.
\end{enumerate}

\subsection{Langkah 3: Layar Login Google (\texttt{lib/screens/login\_screen.dart})}
\begin{enumerate}[leftmargin=*]
    \item \textbf{Navigasi}: Buka folder \texttt{lib/screens/} $\rightarrow$ Buka file \texttt{login\_screen.dart} (Baris 15--42).
    \item \textbf{Fungsi \& Logika}: Membuka modal Google Sign-In dan mengirimkan token OAuth 2.0 ke Firebase Auth.
    \item \textbf{Tampilan Aplikasi}: Menampilkan tombol \textit{''Masuk dengan Google''}. Saat diklik, muncul jendela modal bawaan Android pilihan email Gmail.
\end{enumerate}

\subsection{Langkah 4: Layar Peta Utama (\texttt{lib/screens/home\_screen.dart})}
\begin{enumerate}[leftmargin=*]
    \item \textbf{Navigasi}: Buka folder \texttt{lib/screens/} $\rightarrow$ Buka file \texttt{home\_screen.dart} (Baris 34--39 dan 73--88).
    \item \textbf{Fungsi \& Logika}: Membaca data kenangan Firestore secara \textit{realtime} dan merender peta OpenStreetMap beserta marker pin merah.
    \item \textbf{Tampilan Aplikasi}: Peta interaktif penuh layar yang dapat digeser dan di-zoom, serta memiliki tombol melayang \textit{(+)} di kanan bawah.
\end{enumerate}

\subsection{Langkah 5: Form Tambah Spot (\texttt{lib/screens/add\_spot\_screen.dart})}
\begin{enumerate}[leftmargin=*]
    \item \textbf{Navigasi}: Buka folder \texttt{lib/screens/} $\rightarrow$ Buka file \texttt{add\_spot\_screen.dart} (Baris 32--58, 60--69, dan 84--125).
    \item \textbf{Fungsi \& Logika}: Membaca sensor GPS secara otomatis dan memotret lokasi via kamera HP.
    \item \textbf{Tampilan Aplikasi}: Form input foto/video, kolom nama tempat, catatan, teks koordinat GPS otomatis, dan tombol \textit{''Simpan SpotKu''}.
\end{enumerate}

\subsection{Langkah 6: Rincian Spot \& Demo Notifikasi (\texttt{lib/screens/detail\_screen.dart})}
\begin{enumerate}[leftmargin=*]
    \item \textbf{Navigasi}: Buka \texttt{lib/screens/detail\_screen.dart} dan \texttt{lib/services/notification\_service.dart}.
    \item \textbf{Fungsi \& Logika}: Membuka foto dari memori internal HP dan menguji notifikasi pengingat Android dalam 5 detik.
    \item \textbf{Tampilan Aplikasi}: Foto berukuran besar, teks koordinat desimal, dan tombol pengingat yang memicu \textit{banner} notifikasi melayang di atas layar HP.
\end{enumerate}

---

\section{Naskah Presentasi untuk Dosen Penguji}

Berikut adalah naskah presentasi berdurasi 2--3 menit yang dapat dihafalkan:

\begin{tcolorbox}[colback=teal!5,colframe=teal!60,title=Naskah Presentasi Mahasiswa]
``Selamat pagi/siang Bapak/Ibu Dosen Penguji. Hari ini saya mempresentasikan hasil perancangan project aplikasi Android bernama \textbf{SpotKu}.

\textbf{SpotKu} merupakan aplikasi diari perjalanan digital berbasis peta interaktif yang dibangun menggunakan \textit{framework} Flutter dan \textit{backend} Google Firebase.

Alur teknis utama aplikasi ini mencakup 5 komponen penting:
\begin{enumerate}
    \item \textbf{Autentikasi Aman}: Menggunakan Firebase Auth berbasis Google Sign-In dengan protokol OAuth 2.0.
    \item \textbf{Peta Interaktif Open-Source}: Mengintegrasikan OpenStreetMap melalui paket \texttt{flutter\_map} sehingga aplikasi bebas biaya dan tanpa memerlukan pendaftaran API Key berbayar.
    \item \textbf{Sinkronisasi Realtime}: Memanfaatkan \texttt{StreamBuilder} untuk mendengarkan perubahan koleksi Cloud Firestore secara langsung tanpa perlu \textit{refresh} manual.
    \item \textbf{Integrasi Hardware HP}: Menggunakan \texttt{Geolocator} untuk mendeteksi koordinat GPS otomatis dan \texttt{ImagePicker} untuk mengambil foto/video langsung dari kamera HP.
    \item \textbf{Service Notifikasi}: Menyediakan \texttt{NotificationService} lokal untuk pemicu pengingat kunjungan kembali.
\end{enumerate}

Demikian ringkasan sistem aplikasi SpotKu. Terima kasih dan saya siap menerima masukan atau pertanyaan dari Bapak/Ibu Dosen.''
\end{tcolorbox}

\end{document}
